\documentclass{article}
\usepackage{amsmath,amssymb}
\usepackage{xurl}
\usepackage[hidelinks]{hyperref}
\setlength{\emergencystretch}{2em}
% Shared commands for the notes in docs/paper-gaps/.
% Every note loads this file via % Shared commands for the notes in docs/paper-gaps/.
% Every note loads this file via % Shared commands for the notes in docs/paper-gaps/.
% Every note loads this file via \input{command}. Project-wide notation
% belongs here, never in per-note redefinitions.

% --- Project configuration ------------------------------------------------
\newcommand{\projectIssueBase}{https://github.com/LionSR/QICLean/issues/}
\newcommand{\projectPrBase}{https://github.com/LionSR/QICLean/pull/}
\newcommand{\projectDocBase}{https://sirui-lu.com/QICLean/docs/}
% ---------------------------------------------------------------------------

\DeclareUnicodeCharacter{2080}{\ifmmode{}_{0}\else\textsubscript{0}\fi}
\DeclareUnicodeCharacter{2081}{\ifmmode{}_{1}\else\textsubscript{1}\fi}
\DeclareUnicodeCharacter{2082}{\ifmmode{}_{2}\else\textsubscript{2}\fi}
\DeclareUnicodeCharacter{2099}{\ifmmode{}_{n}\else\textsubscript{n}\fi}

\newcommand{\Lean}{\textsc{Lean}}

% A formal declaration name, upright sans, breakable at underscores.
\ExplSyntaxOn
\NewDocumentCommand{\leanid}{m}{
  \ifmmode
    \textsf{\detokenize{#1}}
  \else
    \group_begin:
    \urlstyle{sf}
    \tl_set:Nn \l_tmpa_tl {#1}
    \tl_replace_all:Nnn \l_tmpa_tl {\_} {_}
    \exp_args:NV \path \l_tmpa_tl
    \group_end:
  \fi
}
\ExplSyntaxOff

% Traceability links; use in footnotes, not in the running prose.
\newcommand{\ghissue}[1]{\href{\projectIssueBase#1}{GitHub issue \##1}}
\newcommand{\ghpr}[1]{\href{\projectPrBase#1}{PR \##1}}
\newcommand{\doclink}[1]{\href{\projectDocBase#1.html}{\path{#1}}}

% Domain notation, using \providecommand so a note-local refinement stays
% possible.
\providecommand{\tr}{\operatorname{tr}}

% Dirac notation and the identity operator, as in blueprint/src/macros.
\usepackage{dsfont}
\providecommand{\ket}[1]{|#1\rangle}
\providecommand{\bra}[1]{\langle #1|}
\providecommand{\braket}[2]{\langle #1 | #2 \rangle}
\providecommand{\ketbra}[2]{|#1\rangle\!\langle #2|}
\providecommand{\Id}{\mathds{1}}
\providecommand{\one}{\mathds{1}}
\providecommand{\C}{\mathbb{C}}
\providecommand{\MN}[1]{M_{#1}(\C)}
\providecommand{\MD}{M_D(\C)}

% Verdict marker: \gapnote{<kind>}{<status>}. Kind is the policy
% classification (clarification, local-correction, scope-restriction,
% unfaithful, false-source, open-gap); status is open, wip (elimination
% underway), resolved, or historical. Machine-read by the site index;
% prints a verdict line.
\newcommand{\gapnote}[2]{\par\noindent\textbf{Verdict:} #1 (#2).\par}
. Project-wide notation
% belongs here, never in per-note redefinitions.

% --- Project configuration ------------------------------------------------
\newcommand{\projectIssueBase}{https://github.com/LionSR/QICLean/issues/}
\newcommand{\projectPrBase}{https://github.com/LionSR/QICLean/pull/}
\newcommand{\projectDocBase}{https://sirui-lu.com/QICLean/docs/}
% ---------------------------------------------------------------------------

\DeclareUnicodeCharacter{2080}{\ifmmode{}_{0}\else\textsubscript{0}\fi}
\DeclareUnicodeCharacter{2081}{\ifmmode{}_{1}\else\textsubscript{1}\fi}
\DeclareUnicodeCharacter{2082}{\ifmmode{}_{2}\else\textsubscript{2}\fi}
\DeclareUnicodeCharacter{2099}{\ifmmode{}_{n}\else\textsubscript{n}\fi}

\newcommand{\Lean}{\textsc{Lean}}

% A formal declaration name, upright sans, breakable at underscores.
\ExplSyntaxOn
\NewDocumentCommand{\leanid}{m}{
  \ifmmode
    \textsf{\detokenize{#1}}
  \else
    \group_begin:
    \urlstyle{sf}
    \tl_set:Nn \l_tmpa_tl {#1}
    \tl_replace_all:Nnn \l_tmpa_tl {\_} {_}
    \exp_args:NV \path \l_tmpa_tl
    \group_end:
  \fi
}
\ExplSyntaxOff

% Traceability links; use in footnotes, not in the running prose.
\newcommand{\ghissue}[1]{\href{\projectIssueBase#1}{GitHub issue \##1}}
\newcommand{\ghpr}[1]{\href{\projectPrBase#1}{PR \##1}}
\newcommand{\doclink}[1]{\href{\projectDocBase#1.html}{\path{#1}}}

% Domain notation, using \providecommand so a note-local refinement stays
% possible.
\providecommand{\tr}{\operatorname{tr}}

% Dirac notation and the identity operator, as in blueprint/src/macros.
\usepackage{dsfont}
\providecommand{\ket}[1]{|#1\rangle}
\providecommand{\bra}[1]{\langle #1|}
\providecommand{\braket}[2]{\langle #1 | #2 \rangle}
\providecommand{\ketbra}[2]{|#1\rangle\!\langle #2|}
\providecommand{\Id}{\mathds{1}}
\providecommand{\one}{\mathds{1}}
\providecommand{\C}{\mathbb{C}}
\providecommand{\MN}[1]{M_{#1}(\C)}
\providecommand{\MD}{M_D(\C)}

% Verdict marker: \gapnote{<kind>}{<status>}. Kind is the policy
% classification (clarification, local-correction, scope-restriction,
% unfaithful, false-source, open-gap); status is open, wip (elimination
% underway), resolved, or historical. Machine-read by the site index;
% prints a verdict line.
\newcommand{\gapnote}[2]{\par\noindent\textbf{Verdict:} #1 (#2).\par}
. Project-wide notation
% belongs here, never in per-note redefinitions.

% --- Project configuration ------------------------------------------------
\newcommand{\projectIssueBase}{https://github.com/LionSR/QICLean/issues/}
\newcommand{\projectPrBase}{https://github.com/LionSR/QICLean/pull/}
\newcommand{\projectDocBase}{https://sirui-lu.com/QICLean/docs/}
% ---------------------------------------------------------------------------

\DeclareUnicodeCharacter{2080}{\ifmmode{}_{0}\else\textsubscript{0}\fi}
\DeclareUnicodeCharacter{2081}{\ifmmode{}_{1}\else\textsubscript{1}\fi}
\DeclareUnicodeCharacter{2082}{\ifmmode{}_{2}\else\textsubscript{2}\fi}
\DeclareUnicodeCharacter{2099}{\ifmmode{}_{n}\else\textsubscript{n}\fi}

\newcommand{\Lean}{\textsc{Lean}}

% A formal declaration name, upright sans, breakable at underscores.
\ExplSyntaxOn
\NewDocumentCommand{\leanid}{m}{
  \ifmmode
    \textsf{\detokenize{#1}}
  \else
    \group_begin:
    \urlstyle{sf}
    \tl_set:Nn \l_tmpa_tl {#1}
    \tl_replace_all:Nnn \l_tmpa_tl {\_} {_}
    \exp_args:NV \path \l_tmpa_tl
    \group_end:
  \fi
}
\ExplSyntaxOff

% Traceability links; use in footnotes, not in the running prose.
\newcommand{\ghissue}[1]{\href{\projectIssueBase#1}{GitHub issue \##1}}
\newcommand{\ghpr}[1]{\href{\projectPrBase#1}{PR \##1}}
\newcommand{\doclink}[1]{\href{\projectDocBase#1.html}{\path{#1}}}

% Domain notation, using \providecommand so a note-local refinement stays
% possible.
\providecommand{\tr}{\operatorname{tr}}

% Dirac notation and the identity operator, as in blueprint/src/macros.
\usepackage{dsfont}
\providecommand{\ket}[1]{|#1\rangle}
\providecommand{\bra}[1]{\langle #1|}
\providecommand{\braket}[2]{\langle #1 | #2 \rangle}
\providecommand{\ketbra}[2]{|#1\rangle\!\langle #2|}
\providecommand{\Id}{\mathds{1}}
\providecommand{\one}{\mathds{1}}
\providecommand{\C}{\mathbb{C}}
\providecommand{\MN}[1]{M_{#1}(\C)}
\providecommand{\MD}{M_D(\C)}

% Verdict marker: \gapnote{<kind>}{<status>}. Kind is the policy
% classification (clarification, local-correction, scope-restriction,
% unfaithful, false-source, open-gap); status is open, wip (elimination
% underway), resolved, or historical. Machine-read by the site index;
% prints a verdict line.
\newcommand{\gapnote}[2]{\par\noindent\textbf{Verdict:} #1 (#2).\par}

\title{The Chebyshev filter near the normalization point}
\author{}
\date{2026-11-04}
\begin{document}
\maketitle
\gapnote{false-source}{resolved}

\section{The printed slope}

Let $\eta_0<\eta_1<u$, put $\gamma=\eta_1-\eta_0$, and let $T_k$ be the
Chebyshev polynomial of the first kind of degree $k\ge1$.
Definition \texttt{def:Cheby-AGSP} of
Arad--Landau--Vazirani--Vidick~\cite{Arad2016RigorousRG}, source lines
795--801, uses $u=\|H\|$ and
\begin{align}
    \tilde P_k(x) & =T_k\left(\frac{2(x-\eta_1)}{u-\eta_1}-1\right),
    \label{eq:alvv_slope_unnormalized}\\
    P_k(x)        & =\frac{\tilde P_k(x)}{\tilde P_k(\eta_0)}.
    \label{eq:alvv_slope_normalized}
\end{align}
Item~2 of Theorem \texttt{thm:Cheby-AGSP} of~\cite{Arad2016RigorousRG},
source lines 822--824, asserts that for $\lambda\le\eta_0+\gamma/k$,
\begin{align}
    P_k(\lambda) & \ge1-O\left(\frac{k|\lambda-\eta_0|}{\gamma u}
    \sqrt\Delta\right),
    \qquad
    \sqrt\Delta=2\exp\left(-2k\sqrt{\frac{\gamma}{u-\eta_0}}\right).
    \label{eq:alvv_slope_printed}
\end{align}
The proof, source lines 844--846, bounds $|T_k(x+\delta)-T_k(x)|$ by
$O(k\delta)$ times a quantity independent of $k$. This omits the growth of
$T_k'$ on $[1,\infty)$, which is exponential in $k$, and the factor $\sqrt\Delta$
in (\ref{eq:alvv_slope_printed}) comes from that omission.

\section{The printed slope fails}

Put $f(x)=(u+\eta_1-2x)/(u-\eta_1)$. For $x\le\eta_1$ the argument in
(\ref{eq:alvv_slope_unnormalized}) is $-f(x)\le-1$, and the parity
$T_k(-y)=(-1)^kT_k(y)$ gives $P_k(x)=T_k(f(x))/T_k(f(\eta_0))$.
Write $f(x)=\cosh s$ and $f(\eta_0)=\cosh s_0$ with $0\le s\le s_0$.
Since $T_k(\cosh t)=\cosh(kt)$,
\begin{align}
    P_k(x) & =\frac{\cosh(ks)}{\cosh(ks_0)}\le2e^{-k(s_0-s)}.
    \label{eq:alvv_slope_cosh_ratio}
\end{align}
Convexity of $\cosh$ gives $\cosh s_0-\cosh s\le\sinh s_0\,(s_0-s)$. At
$x=\eta_0+\gamma/k$,
\begin{align}
    \cosh s_0-\cosh s & =\frac{2\gamma}{k(u-\eta_1)},
    \label{eq:alvv_slope_gap}\\
    \sinh s_0         & =\sqrt{f(\eta_0)^2-1}
    =\frac{2\sqrt{\gamma(u-\eta_0)}}{u-\eta_1}.
    \label{eq:alvv_slope_sinh}
\end{align}
Hence $k(s_0-s)\ge\sqrt{\gamma/(u-\eta_0)}$, and
(\ref{eq:alvv_slope_cosh_ratio}) yields
\begin{align}
    P_k\left(\eta_0+\frac\gamma k\right)
     & \le2\exp\left(-\sqrt{\frac{\gamma}{u-\eta_0}}\right).
    \label{eq:alvv_slope_upper}
\end{align}
The right side of (\ref{eq:alvv_slope_upper}) does not depend on $k$.
For $\eta_0=0$, $\eta_1=1$ and $u=2$ it equals $2e^{-1/\sqrt2}<0.99$. At
$\lambda=\eta_0+\gamma/k$ the right side of (\ref{eq:alvv_slope_printed}) is
$1-O(e^{-\sqrt2\,k})$, which exceeds $0.99$ for all large $k$, whatever the
implied constant. The printed bound therefore fails. Numerically,
$P_{20}(1/20)=0.486\ldots$ for these parameters.\footnote{Formal declaration:
    \leanid{Polynomial.Chebyshev.eval_T_div_eval_T_le_two_mul_exp}, which
    proves (\ref{eq:alvv_slope_upper}) for the explicit polynomial
    (\ref{eq:alvv_slope_normalized}).}

\section{The corrected slope}

For $1\le y\le y_0$, write $y=\cosh s$, $y_0=\cosh t$ with $0\le s\le t$.
The ratio $\sinh(ks)/\sinh s$ is nondecreasing in $s>0$, by induction on $k$
from $\sinh((k+1)s)=\sinh(ks)\cosh s+\cosh(ks)\sinh s$. Therefore
$x\mapsto k\sinh(kt)\cosh x-\sinh t\cosh(kx)$ is nondecreasing on $[0,t]$,
and evaluating it at $s$ and $t$ gives
\begin{align}
    \sqrt{y_0^2-1}\,\bigl(T_k(y_0)-T_k(y)\bigr)
     & \le k(y_0-y)\,T_k(y_0).
    \label{eq:alvv_slope_tangent}
\end{align}
Here $\sinh(kt)\le\cosh(kt)=T_k(y_0)$ was used. Apply
(\ref{eq:alvv_slope_tangent}) with $y=f(\lambda)$ and $y_0=f(\eta_0)$ for
$\eta_0\le\lambda\le\eta_1$. Then $y_0-y=2(\lambda-\eta_0)/(u-\eta_1)$ and
(\ref{eq:alvv_slope_sinh}) give, for every $\lambda\le\eta_1$,
\begin{align}
    P_k(\lambda) & \ge1-\frac{k|\lambda-\eta_0|}{\sqrt{\gamma(u-\eta_0)}}.
    \label{eq:alvv_slope_corrected}
\end{align}
For $\lambda\le\eta_0$ this is weaker than $P_k(\lambda)\ge1$. The bound
(\ref{eq:alvv_slope_corrected}) holds on all of $\lambda\le\eta_1$, which
contains the printed range $\lambda\le\eta_0+\gamma/k$; there the loss is at
most $\sqrt{\gamma/(u-\eta_0)}$. It replaces (\ref{eq:alvv_slope_printed}):
the factor $\sqrt\Delta$ is removed and $\gamma u$ is replaced by
$\sqrt{\gamma(u-\eta_0)}$. Its order in $k$ and $|\lambda-\eta_0|$ is sharp:
the derivative of $P_k$ at $\eta_0$ is
$-k\tanh(ks_0)/\sqrt{\gamma(u-\eta_0)}$.\footnote{Formal declarations:
    \leanid{Polynomial.Chebyshev.sqrt_mul_eval_T_sub_le} for
    (\ref{eq:alvv_slope_tangent}),
    \leanid{Polynomial.Chebyshev.exists_interval_shrinking_polynomial_one_sub_le}
    for (\ref{eq:alvv_slope_corrected}), and
    \leanid{Submodule.exists_isSpectralProjector_chebyshev_one_sub_le} for
    the operator form.}

\section{The rank base in item 4}

Item~4 of the same theorem, source lines 825--826, bounds the Schmidt rank
of $K=P_k(H)$ at a cut inside the middle region by $\tilde B^{O(k)}$, where
$\tilde B$ bounds the Schmidt rank of the middle term $H_M$. At such a cut,
$H=A\otimes1+H_M+1\otimes C$, where $A$ collects the terms left of the cut
and $C$ those to its right. Hence $H$ has Schmidt rank at most $\tilde B+2$,
and $K$, a combination of the powers $H^j$ with $j\le k$, has Schmidt rank at
most
\begin{align}
    \sum_{j=0}^k(\tilde B+2)^j & \le(\tilde B+2)^{k+1}.
    \label{eq:alvv_slope_rank}
\end{align}
For $\tilde B\ge2$ this is at most $\tilde B^{2k+2}$, in the printed form.
For $\tilde B\le1$ the base $\tilde B$ cannot absorb the two boundary
terms. If $H_M=0$, $A=\operatorname{diag}(\alpha_i)$ and
$C=\operatorname{diag}(\beta_j)$ with distinct entries, then the Schmidt rank
of $K$ is the rank of the matrix $\bigl(P_k(\alpha_i+\beta_j)\bigr)_{ij}$,
which is $k+1$ for suitable entries, not $1$. The bound
(\ref{eq:alvv_slope_rank}) is the one formalized.
\end{document}
